\usepackage{ifluatex}

\usepackage{amssymb,amsthm}
\usepackage[fleqn]{amsmath}
\usepackage[svgnames]{xcolor}
\usepackage{graphicx}
\usepackage{hyperref}
\hypersetup{
	pdftitle={Math 121B},
	%pdfauthor={Neil Donaldson},
}
\usepackage[breakable,theorems,skins,hooks]{tcolorbox}

%lualatex
\ifluatex
%\usepackage{fontspec}
%\usepackage{unicode-math}
\usepackage{microtype}
\usepackage[otfmath,theoremfont,trueslanted,largesc]{newpx}
\setmathfont{TeX Gyre Pagella Math}
\setmathfont{TeX Gyre Pagella Math}[range="222B-"2233,Scale=1.08]
\setmathfont{TeX Gyre Pagella Math}[range=\sum]
\everydisplay{\Umathoperatorsize\displaystyle=4.5ex}

%pdflatex
\else
% %\usepackage[utf8]{inputenc}
\usepackage[T1]{fontenc}
\usepackage[theoremfont,trueslanted,largesc]{newpx}
\fi


%lengths
\unitlength 1cm
\setlength{\textheight}{22.5cm}
%Reset from 0pt plus 1pt
\setlength{\parskip}{0pt plus 2pt minus 0.25pt}
\parindent0pt
\textwidth 17cm
\oddsidemargin -0.5cm
\evensidemargin -0.5cm
\topmargin -1.5cm
\topskip 0cm
\headheight 0.5cm
\headsep 1cm
\marginparwidth 1.2cm

\newlength\doubleind
\addtolength{\doubleind}{\leftmargini}
\addtolength{\doubleind}{\leftmarginii}

\newlength\doubleindentmini
\addtolength\doubleindentmini{\textwidth}
\addtolength\doubleindentmini{-\doubleind}


%MathNumbering
\makeatletter
\@addtoreset{equation}{section}
\renewcommand{\theequation}{\thesection.\@arabic\c@equation}
\newcommand\Nopagebreak{\@nobreaktrue\nopagebreak}
%\renewcommand{\theequation}{\@arabic\c@equation}
\makeatother

\newcommand\boldinline[1]{\paragraph{#1}}
\newcommand\boldsubsection[1]{\subsection*{#1}}
\newcommand\boldsubsubsection[1]{\subsubsection*{#1}}

\def\exstart{\hangindent\leftmargini\textup{1.}\hspace{\labelsep}}

\makeatletter
\newenvironment{enumeratepic}{
  \@mplistdepth=1
	\begin{enumerate}
}{
	\end{enumerate}
}
\newenvironment{itemizepic}{
	\makeatletter
  \@mplistdepth=1
  \makeatother
	\begin{itemize}
}{
	\end{itemize}
}
\makeatother
             
\newenvironment{enumeratea}{
	\begin{enumerate}
	  \renewcommand{\labelenumi}{(\alph{enumi})}}
	{\end{enumerate}}


%Tcolorbox Theorems
\tcbset{
	defstyle/.style={enhanced, top=3pt, bottom=3pt, colframe=black, coltitle=black, arc=5pt, boxrule=1.5pt, left*=0pt, right*=0pt, theorem style=plain, terminator sign={.\ \ \ }, fonttitle=\bfseries\upshape, fontupper=\upshape, colback=blue!7!white, grow sidewards by=8pt, drop fuzzy shadow},
	exstyle/.style={enhanced, breakable, beforeafter skip balanced=10pt, coltitle=black, theorem style=plain, terminator sign={.\ \ \ }, fonttitle=\bfseries\upshape, fontupper=\upshape, blanker, borderline west={4pt}{-8pt}{orange!75!white}},
	thmstyle/.style={enhanced, top=3pt, bottom=3pt, colframe=black, coltitle=black, arc=5pt, boxrule=1.5pt,left*=0pt, right*=0pt, theorem style=plain, terminator sign={.\ \ \ }, fonttitle=\bfseries\upshape, fontupper=\slshape, colback=green!12!white, grow sidewards by=8pt, drop fuzzy shadow},
	exercisestyle/.style={enhanced, breakable, beforeafter skip balanced=10pt, coltitle=black, theorem style=plain, terminator sign={.\ \ \ }, fonttitle=\bfseries\upshape, fontupper=\upshape, blanker, borderline west={4pt}{-8pt}{purple!75!white}, title={Exercises\ \ \ }},
	proofstyle/.style={enhanced, breakable, beforeafter skip balanced=10pt, blanker, borderline west={4pt}{-8pt}{green!50!white}},
	asidestyle/.style={enhanced, breakable, beforeafter skip balanced=10pt, blanker, borderline west={4pt}{-8pt}{black!50!white}},
	exercisestyle2/.style={enhanced, breakable, beforeafter skip balanced=10pt, coltitle=black, theorem style=plain, terminator sign={.\ \ \ }, fonttitle=\bfseries\upshape, fontupper=\upshape, blanker, borderline west={4pt}{-8pt}{purple!75!white}, title={Exercises\ \thesubsection.\ \ \ }},
	exercisestyle4/.style={enhanced, breakable, beforeafter skip balanced=10pt, coltitle=black, theorem style=plain, terminator sign={.\ \ \ }, fonttitle=\bfseries\upshape, fontupper=\upshape, blanker, borderline west={4pt}{-8pt}{purple!75!white}, title={Exercises\ \thesection.\ \ \ }},
	exercisestyle3/.style={enhanced, breakable, beforeafter skip balanced=10pt, coltitle=black, theorem style=plain, terminator sign={.\ \ \ }, fonttitle=\bfseries\upshape, fontupper=\upshape, blanker, borderline west={4pt}{-8pt}{purple!75!white}, title={Exercise\quad}},
	conjstyle/.style={enhanced, top=3pt, bottom=3pt, colframe=black, coltitle=black, arc=5pt, boxrule=1.5pt,left*=0pt, right*=0pt, theorem style=plain, terminator sign={.\ \ \ }, fonttitle=\bfseries\upshape, fontupper=\slshape, colback=red!12!white, grow sidewards by=8pt, drop fuzzy shadow},
	}

\tcolorboxenvironment{proof}{breakable,proofstyle}

\newtcbtheorem[number within=section]{defn}{Definition}{defstyle}{defn}
\newtcbtheorem[use counter from=defn]{example}{Example}{exstyle}{ex}
\newtcbtheorem[use counter from=defn]{examples}{Examples}{exstyle}{ex}
\newtcbtheorem[use counter from=defn]{thm}{Theorem}{thmstyle}{thm}
\newtcbtheorem[use counter from=defn]{lemm}{Lemma}{thmstyle}{lemm}
\newtcbtheorem[use counter from=defn]{cor}{Corollary}{thmstyle}{cor}
\newtcbtheorem[use counter from=defn]{axiom}{Axiom}{defstyle}{axiom}
\newtcbtheorem[use counter from=defn]{axioms}{Axioms}{defstyle}{axioms}
\newtcbtheorem[use counter from=defn]{conj}{Conjecture}{conjstyle}{conj}

\newtcolorbox{aside}{asidestyle}
\newtcolorbox{exercises*}{exercisestyle}
\newtcolorbox{exercises}{exercisestyle2}
\newtcolorbox{exercisessec}{exercisestyle4}
\newtcolorbox{exercise}{exercisestyle3}
